\section{Test 2: radicals i potències}
\iniciTest{r2}{b,d,b,c,a,c,d,b,c,a}

\begin{enumerate}
\pregunta Quina de les següents afirmacions és vertadera?

\begin{enumerate}
\opcio{a}{L’arrel $n$-èsima d’un nombre negatiu sempre existeix
}
\opcio{b}{L’arrel $\sqrt{x^{2}}$ té per valor $\left\vert x\right\vert $
}
\opcio{c}{L’equació $x^{4}=-3^{4}$ té per solució $-3$
}
\opcio{d}{$\sqrt{a^{2}+b^{2}}=a+b$, qualssevol que siguin els nombres reals
$a$ i $b$}
\end{enumerate}
\feedback


\pregunta Quina de les següents afirmacions és falsa?

\begin{enumerate}
\opcio{a}{$2^{-3/5}=1/\sqrt[5]{2^{3}}$
}
\opcio{b}{$\sqrt[3]{2^{5}}/2=2^{2/3}$
}
\opcio{c}{$\sqrt{2\sqrt{2\sqrt{2}}}=2^{7/8}$
}
\opcio{d}{$3\cdot 32^{4/5}=16$}
\end{enumerate}
\feedback


\pregunta En convertir en una sola arrel l’expressió%
\[
\frac{\sqrt[4]{a^{3}}\cdot \sqrt[3]{a}\cdot \sqrt{a}}{a\cdot \sqrt[6]{a}}
\]%
s’obté:

\begin{enumerate}
\opcio{a}{$\sqrt{a}$
}
\opcio{b}{$\sqrt[12]{a^{5}}$
}
\opcio{c}{$\sqrt[3]{a}$
}
\opcio{d}{Cap de les anteriors}
\end{enumerate}
\feedback


\pregunta En simplificar l’expressió%
\[
\sqrt[3]{16}-\sqrt[3]{81}+\sqrt[3]{250}+\sqrt[3]{\frac{3}{8}}
\]%
s’obté:

\begin{enumerate}
\opcio{a}{$\frac{9}{2}\cdot \sqrt[3]{5}$
}
\opcio{b}{$7\cdot \sqrt[3]{3}$
}
\opcio{c}{$7\cdot \sqrt[3]{2}-\frac{5}{2}\sqrt[3]{3}$
}
\opcio{d}{Cap de les anteriors}
\end{enumerate}
\feedback


\pregunta En calcular el valor de l’expressió%
\[
(3-\sqrt{6})\cdot (\sqrt{3}+\sqrt{2})-(2\sqrt{3}-3\sqrt{2})^{2}+3\sqrt{6}
\]%
s’obté:

\begin{enumerate}
\opcio{a}{$-30+\sqrt{3}+15\sqrt{6}$
}
\opcio{b}{$-14\sqrt{6}$
}
\opcio{c}{\thinspace $-45$
}
\opcio{d}{Cap de les anteriors}
\end{enumerate}
\feedback


\pregunta En convertir en una sola potència de $a$ l’expressió%
\[
\sqrt[4]{\frac{a^{2}\cdot \sqrt{a}}{\sqrt[4]{a}\cdot \left( \sqrt[3]{a}%
\right) ^{7}}}
\]%
s’obté:

\begin{enumerate}
\opcio{a}{$1$
}
\opcio{b}{$a^{-1/12}$
}
\opcio{c}{$a^{-1/48}$
}
\opcio{d}{$a$}
\end{enumerate}
\feedback


\pregunta Quina de les següents afirmacions és falsa?

\begin{enumerate}
\opcio{a}{$\sqrt[5]{4}\leq \sqrt[10]{17}\leq \sqrt[6]{6}$
}
\opcio{b}{$\sqrt[3]{4}=\sqrt[6]{16}=\sqrt[9]{64}$
}
\opcio{c}{$\frac{2}{\sqrt[3]{2}}=\sqrt[3]{2}=\frac{2}{\sqrt[6]{4}}$
}
\opcio{d}{$\frac{1}{\sqrt[3]{2}}\leq \frac{1}{\sqrt[4]{3}}\leq \frac{1}{\sqrt[6]{%
5}}$}
\end{enumerate}
\feedback


\pregunta En racionalitzar l’expressió
\[
\frac{3\sqrt{3}}{\sqrt[3]{16}}
\]%
s’obté:

\begin{enumerate}
\opcio{a}{$\frac{3}{4}\sqrt[3]{12}$
}
\opcio{b}{$\frac{3}{4}\sqrt[6]{432}$
}
\opcio{c}{$\frac{3}{2}\sqrt[3]{6}$
}
\opcio{d}{$\frac{3}{4}\sqrt{12}$}
\end{enumerate}
\feedback


\pregunta En racionalitzar l’expressió%
\[
\frac{4\sqrt{2}}{5\sqrt{3}-2\sqrt{2}}
\]%
s’obté:

\begin{enumerate}
\opcio{a}{$\frac{36\sqrt{6}}{83}$
}
\opcio{b}{$\frac{36\sqrt{6}}{67}$
}
\opcio{c}{$\frac{16+20\sqrt{6}}{67}$
}
\opcio{d}{Cap de les anteriors}
\end{enumerate}
\feedback


\pregunta En simplificar l’expressió%
\[
\sqrt{a\cdot \sqrt[3]{a^{2}}}\cdot \sqrt[3]{a\cdot \sqrt{a^{3}}}\cdot \sqrt[3%
]{a\cdot \sqrt[4]{a^{3}}}\cdot \sqrt[3]{a^{2}\cdot \sqrt{a\cdot \sqrt{a}}}
\]%
s’obté

\begin{enumerate}
\opcio{a}{$a^{3}\cdot \sqrt[6]{a}$
}
\opcio{b}{$a^{2}\cdot \sqrt[3]{a}$
}
\opcio{c}{$\sqrt[6]{a}$
}
\opcio{d}{Cap de les anteriors}
\end{enumerate}
\feedback

\end{enumerate}